
\clearpage
\VerdierPhysicalPageAnchor{0130}%
\begingroup
\small
\setlength{\emergencystretch}{1.5em}%
\setlength{\abovedisplayskip}{2pt plus 1pt minus 1pt}%
\setlength{\belowdisplayskip}{2pt plus 1pt minus 1pt}%

\VerdierSemanticAnchor{ch2-2-1-3-2-sm2-completion}
\noindent を次の可換図式に補完できることを示せばよい：

\VerdierSemanticAnchor{ch2-2-1-3-2-completion-diagram}
\begin{center}
  $t\in S$\qquad
  \begin{tikzcd}[column sep=4.0em,row sep=2.5em,baseline=(current bounding box.center)]
    Z' \arrow[r,dashed,"g"] \arrow[d,dashed,"t"'] &
    X' \arrow[d,"s"] \\
    Z \arrow[r,"f"'] & X
  \end{tikzcd}
  \quad ,
\end{center}
\VerdierSemanticAnchor{ch2-2-1-3-2-opposite-reduction}
\noindent そうすれば、SM2) の第二の主張は反対圏へ移ることで従う。

\VerdierSemanticAnchor{ch2-2-1-3-2-triangle-construction}
TRI) と TRII) により完全三角 $(X,Y,Z[1],u,v,f[1])$ が存在し、TRI)
により完全三角 $(X',Y,Z',us,v',w')$ が存在する。よって次の可換図式を得る：

\VerdierSemanticAnchor{ch2-2-1-3-2-us-square}
\begin{center}
  \begin{tikzcd}[column sep=4.2em,row sep=2.5em]
    X' \arrow[r,"us"] \arrow[d,"s"'] &
    Y \arrow[d,"\operatorname{id}_Y"] \\
    X \arrow[r,"u"'] & Y
  \end{tikzcd}
  \quad .
\end{center}
\VerdierSemanticAnchor{ch2-2-1-3-2-sm6-application}
\noindent 射 $\operatorname{id}_Y$ は $S$ に属し (SM1)、SM6)
により、$(s,\operatorname{id}_Y,s')$ が三角の射となるような $S$ に属する
射 $s'\colon Z'\to Z[1]$ を得る。特に、次の図式は可換である。

\VerdierSemanticAnchor{ch2-2-1-3-2-shift-square}
\begin{center}
  \begin{tikzcd}[column sep=4.2em,row sep=2.5em]
    Z' \arrow[r,"{w'}"] \arrow[d,"{s'}"'] &
    X'[1] \arrow[d,"{s[1]}"] \\
    Z[1] \arrow[r,"{f[1]}"'] & X[1]
  \end{tikzcd}
  \quad ,
\end{center}

\VerdierSemanticAnchor{ch2-2-1-3-2-proof-complete}
\noindent したがって SM5) により SM2) の第一の主張が従う。

\VerdierSemanticAnchor{ch2-2-1-4}
\medskip
\noindent\textbf{2.1.4.}\quad
\VerdierSemanticAnchor{ch2-2-1-4-convention}
\emph{特に反対のことを明記しない限り、三角圏の乗法的系は三角構造と
両立するものとする。したがって、言葉を濫用して単に乗法的系と呼ぶ。}

\VerdierSemanticAnchor{ch2-2-1-5}
\medskip
\noindent\textbf{2.1.5.}\quad
\VerdierSemanticAnchor{ch2-2-1-5-full-triangulated-subcategory}
$\mathcal D$ を三角圏とする。$\mathcal D$ の\emph{充満三角部分圏}とは、
シフトで安定な $\mathcal D$ の充満部分圏であって、包含関手が（シフトと
厳密に可換するものとして）完全関手となるような三角構造を備えたものをいう。

\VerdierSemanticAnchor{ch2-2-1-5-additive-subcategory}
$\mathcal D$ の充満三角部分圏は $\mathcal D$ の加法部分圏である (1.2.7)。

\VerdierSemanticAnchor{ch2-2-1-5-boundary}
$\mathcal D$ の充満部分圏
\endgroup
